\documentclass{article}
\usepackage[T1]{fontenc}
\usepackage[utf8]{inputenc}
\usepackage{amsmath,amssymb,amsthm,booktabs,hyperref,bookmark}
\usepackage[margin=1in]{geometry}
\usepackage{longtable,listings,xcolor}
\usepackage{graphicx}
\lstset{extendedchars=false,basicstyle=\ttfamily\scriptsize,breaklines=true,columns=fullflexible,keepspaces=true,showstringspaces=false,frame=single,rulecolor=\color{gray!40},numbers=left,numberstyle=\tiny\color{gray},numbersep=4pt}
\lstdefinestyle{python}{language=Python,keywordstyle=\color{blue!60!black}\bfseries,commentstyle=\color{green!40!black}\itshape,stringstyle=\color{red!60!black}}
\lstdefinestyle{json}{language={},morestring=[b]",morestring=[b]',sensitive=true}
\lstdefinestyle{text}{language={}}
\lstdefinestyle{texfrag}{language=[LaTeX]TeX,basicstyle=\ttfamily\scriptsize}
\emergencystretch=12em
\sloppy
\allowdisplaybreaks

\newtheorem{proposition}{Proposition}
\newtheorem{theorem}[proposition]{Theorem}
\newtheorem{lemma}[proposition]{Lemma}
\newtheorem{corollary}[proposition]{Corollary}
\newtheorem{definition}[proposition]{Definition}
\newtheorem{gate}[proposition]{Gate}
\newtheorem{inputstatement}[proposition]{Input}
\theoremstyle{remark}
\newtheorem{remark}[proposition]{Remark}
\newtheorem{nonclaim}[proposition]{Non-claim}

\newcommand{\scopebox}[2]{%
  \par\medskip\noindent\fbox{\begin{minipage}{0.92\textwidth}
  \textbf{#1}\par\smallskip #2\end{minipage}}\par\medskip}
\newcommand{\Pc}{\mathcal P}
\newcommand{\sigeps}{\sigma_\varepsilon}
\newcommand{\chiI}{\chi_{\mathrm I}}
\newcommand{\Ro}{R_0}
\newcommand{\Rp}{R_0^{\mathrm{pred}}}
\newcommand{\Ramp}{\mathfrak R}
\newcommand{\Dsrc}{\mathcal D}
\newcommand{\Ncal}{\mathcal N}
% Extra macros that appear in provenance manuscripts (no-ops / stubs if unused)
\newcommand{\RL}{\mathfrak R_{\mathrm L}}
\newcommand{\RK}{\mathfrak R_{\mathrm K}}
\newcommand{\RM}{\mathfrak R_{\mathrm M}}
\newcommand{\JD}{\mathcal J_{\mathrm D}}
\newcommand{\Ccoh}{\mathcal C_{\mathrm{coh}}}
\newcommand{\Creg}{\mathcal C_{\mathrm{reg}}}
\newcommand{\Tcol}{\mathcal T_{\mathrm{collision}}}
\newcommand{\Coll}{\mathrm{Collapse}}
\newcommand{\EA}{E_0^{\mathrm A}}
\newcommand{\ER}{E_0^{\mathrm R}}
\newcommand{\EAm}{E_m^{\mathrm A}}
\newcommand{\ERm}{E_m^{\mathrm R}}
\newcommand{\PP}{\mathcal P}
\newcommand{\sig}{\sigma_\varepsilon}
\newcommand{\MT}{M_T}
\newcommand{\MTo}{M_T^\circ}
\newcommand{\Pmass}{P_m}
\newcommand{\Bcal}{\mathcal B}
\newcommand{\Isol}{\mathcal I_{\mathrm{sol}}}
\newcommand{\mc}{\mathscr F_{\mathrm{mc,pol}}}
\providecommand{\dd}{\mathrm{d}}
\providecommand{\R}{\mathbb{R}}

\title{Landmark XVI-N\\
\large Classical completion --- PreparedClassicalBVInterfaceThroughOrderKappa\\
\large Combined manuscript (scientific text + scripts + certificates + provenance)}
\author{XVI-N r12 release}
\date{2026-10-06}

\begin{document}
\maketitle
\noindent\textbf{Combined package.} This PDF is the self-contained XVI-N release:
Part~I = scientific manuscript (r12 release; Route B unresolved);
Part~II = all Python scripts and certificates;
Part~III = supporting notes;
Part~IV = provenance predecessors XI-A, XI-B, XII-C, XVI-M and the release manifest.

\tableofcontents
\newpage

%=====================================================================
\part{Scientific manuscript (XVI-N)}
%=====================================================================
% Scientific abstract and body follow (from the r9 XVI-N.tex).
\input{body/XVI-N-body.tex}

%=====================================================================
\part{Computational appendix}
%=====================================================================
\input{body/part2_listings.tex}

%=====================================================================
\part{Notes appendix}
%=====================================================================
\input{body/part3_notes.tex}

%=====================================================================
\part{Provenance and references}
%=====================================================================
% Provenance manuscripts are included as verbatim source listings to avoid
% macro clashes with the main document while preserving the full original text
% and recorded hashes. Cleaned body extracts also ship under provenance/*-body.tex.
\section{Provenance manuscripts (full source listings)}
\label{app:prov}
The XVI-M baseline sha256 is
\texttt{41aa109935f40ddee3109de62c46329c189b5408007970cc90ed5611dc4739ea}
(verified unchanged against the package baseline).

\subsection{Baseline hashes and manifests}
\lstinputlisting[style=text]{listing_src/provenance/baseline_sha256.txt}
\lstinputlisting[style=text]{listing_src/provenance/baseline_sha256_revision.txt}
\lstinputlisting[style=json]{listing_src/provenance/N_dependency_manifest.json}
\lstinputlisting[style=json]{listing_src/provenance/XVI-M_conventions.json}

\subsection{Landmark XI-A (full source)}
\label{app:XIA}
\textbf{sha256:} see inventory / \texttt{sha256sum provenance/XI-A.tex}.
\lstinputlisting[style=texfrag,caption={XI-A.tex}]{listing_src/provenance/XI-A.tex}

\subsection{Landmark XI-B (full source)}
\label{app:XIB}
\lstinputlisting[style=texfrag,caption={XI-B.tex}]{listing_src/provenance/XI-B.tex}

\subsection{Landmark XII-C (full source)}
\label{app:XIIC}
\lstinputlisting[style=texfrag,caption={XII-C.tex}]{listing_src/provenance/XII-C.tex}

\subsection{Landmark XVI-M (full source)}
\label{app:XVIM}
\lstinputlisting[style=texfrag,caption={XVI-M.tex}]{listing_src/provenance/XVI-M.tex}

\subsection{RELEASE\_MANIFEST}
\label{app:manifest}
\lstinputlisting[style=text]{listing_src/provenance/RELEASE_MANIFEST.txt}

\end{document}
